\documentclass[12pt,a4paper]{article}
\usepackage[utf8]{inputenc}
\usepackage[T1]{fontenc}
\usepackage[spanish]{babel}
\usepackage{geometry}
\usepackage{xcolor}
\usepackage{titlesec}
\usepackage{enumitem}
\usepackage{tcolorbox}
\usepackage{fontawesome5}

\geometry{top=1in, bottom=1in, left=1in, right=1in}

% Colors
\definecolor{primary}{HTML}{2C3E50}
\definecolor{secondary}{HTML}{3498DB}
\definecolor{accent}{HTML}{E67E22}
\definecolor{bg}{HTML}{FDFEFE}

% Styles
\titleformat{\section}{\large\bfseries\color{primary}}{}{0em}{\faBook \quad}
\titleformat{\subsection}{\normalsize\bfseries\color{secondary}}{}{0em}{}

\newtcolorbox{note}[1]{
    colback=bg,
    colframe=secondary,
    fonttitle=\bfseries,
    title=#1,
    arc=3pt,
    boxrule=1pt
}

\newtcolorbox{trap}[1]{
    colback=red!5,
    colframe=accent,
    fonttitle=\bfseries,
    title=\faExclamationTriangle \quad #1,
    arc=3pt,
    boxrule=1pt
}

\begin{document}

\begin{center}
    {\Huge \color{primary} \textbf{Apuntes de Inglés}} \\
    \vspace{0.2cm}
    {\large \color{secondary} Dominando de poco en poco}
\end{center}

\vspace{0.5cm}

\section{Pronombres de Objeto (It vs. Them)}

Los pronombres de objeto reciben la acción y van \textbf{después} del verbo o preposición.

\begin{note}{1. La Regla de la Cantidad 🔢}
    \begin{itemize}[noitemsep]
        \item \textbf{It (Singular):} Para una sola cosa o conceptos incontables. \\
              \textit{Ejemplo:} The phone $\to$ \textbf{It}.
        \item \textbf{Them (Plural):} Para varias cosas o personas. \\
              \textit{Ejemplo:} The keys $\to$ \textbf{Them}.
    \end{itemize}
\end{note}

\begin{note}{2. Objetos con "S" que SÍ son Plurales 👓}
    Objetos de dos partes que siempre usan \textbf{Them}:
    \begin{itemize}[noitemsep]
        \item \textit{Glasses} (Lentes), \textit{Scissors} (Tijeras), \textit{Pants/Jeans} (Pantalones), \textit{Pyjamas} (Pijamas).
        \item \textbf{Correcto:} "Where \textbf{are} my glasses? I can't find \textbf{them}."
    \end{itemize}
\end{note}

\begin{trap}{3. La Trampa de la "S" Falsa (Singulares) 🪤}
    Palabras que terminan en "s" pero se tratan como un paquete único. Usan \textbf{It}:
    \begin{itemize}[noitemsep]
        \item \textbf{News} (Noticias): "I love the news, I watch \textbf{it} every day."
        \item \textbf{Materias académicas:} \textit{Physics} (Física), \textit{Mathematics} (Matemáticas), \textit{Economics} (Economía).
        \item \textbf{Juegos:} \textit{Darts} (Dardos), \textit{Billiards} (Billar), \textit{Dominoes} (Dominó).
    \end{itemize}
    \textbf{Truco:} Si no puedes quitarle la "s" para tener uno solo (no existe "a new"), es singular.
\end{trap}

\begin{note}{4. Posición en Phrasal Verbs (El Sándwich) 🥪}
    El pronombre \textbf{siempre} va en medio de las dos palabras del verbo.
    \begin{itemize}[noitemsep]
        \item \textbf{Pick them up} (Recógelas).
        \item \textbf{Turn it down} (Bájale).
    \end{itemize}
\end{note}

\section{Registro de Mis Correcciones Clave}
\begin{itemize}
    \item \textbf{Decisiones al momento:} "I\textbf{'ll} try them on" (Me los probaré).
    \item \textbf{Órdenes negativas:} "\textbf{Don't} throw it out" (No lo tires).
    \item \textbf{Demostrativos:} "\textbf{This} box" (singular) vs "\textbf{These} boxes" (plural).
\end{itemize}

\end{document}
